% =====================================================================
%  M1 -- Stationary points of a quartic landscape
% =====================================================================
\subsection{M1: Stationary points of a quartic landscape}
\setprob{M1}

\begin{frame}{\probtitle{M1}{Stationary points of a quartic landscape}}
  \begin{probbox}[Problem M1 \hfill Mathematical \quad $\star\star$]
    Consider \(\displaystyle f(x,y) = x^4 + y^4 - 4xy\) on \(\R^2\).
    \begin{enumerate}[(a)]
      \item Compute $\nabla f$ and find \emph{all} stationary points.
      \item Classify each stationary point using the Hessian.
      \item Take one gradient-descent step from $\x_0=(1,0)$ with $\alpha = 0.1$,
            and check that $f$ goes down.
      \item Near the minimisers, which fixed step sizes keep GD locally stable?
            What happens if GD starts exactly at the saddle, or anywhere on the line $y=-x$?
    \end{enumerate}
  \end{probbox}
  \footnotesize\textcolor{mGrey}{What this practises: computing gradients, finding and classifying stationary points,
  and local step-size stability.}
\end{frame}

\begin{frame}{\soltitle{M1}{1}{4}}
  \textbf{(a) Gradient and stationary points.}
  \[
    \nabla f(x,y) = \begin{pmatrix} 4x^3 - 4y \\ 4y^3 - 4x \end{pmatrix} = \mathbf 0
    \quad\Longleftrightarrow\quad y = x^3, \;\; x = y^3 .
  \]
  Substitute the first equation into the second:
  \[
    x = (x^3)^3 = x^9 \;\Longrightarrow\; x\,(x^8 - 1) = 0 .
  \]
  Over $\R$, $x^8 = 1$ has only the roots $x = \pm 1$. So $x \in \{-1, 0, 1\}$, and with $y = x^3$ we get three points:
  \[
    \ans{(0,0),\quad (1,1),\quad (-1,-1)}
  \]
\end{frame}

\begin{frame}{\soltitle{M1}{2}{4}}
  \textbf{(b) Classify with the Hessian.}\quad
  $\displaystyle \nabla^2 f = \begin{pmatrix} 12x^2 & -4 \\ -4 & 12y^2 \end{pmatrix}.$
  \vspace{2pt}
  \begin{itemize}\small
    \item At $(0,0)$: $\begin{pmatrix} 0 & -4\\-4&0\end{pmatrix}$ has eigenvalues $-4$ (along $(1,1)$) and $+4$ (along $(1,-1)$).
          Indefinite, so this is a \alert{saddle point}, with $f=0$.
    \item At $(\pm1,\pm1)$: $\begin{pmatrix} 12 & -4\\-4&12\end{pmatrix}$ has eigenvalues $8$ (along $(1,1)$) and $16$ (along $(1,-1)$).
          Positive definite, so these are \alert{strict local minima}, with $f = 1+1-4 = -2$.
    \item Are they global? Since $4|xy| \le 2(x^2+y^2)$, we have $f \ge x^4+y^4-2(x^2+y^2) \to \infty$, so $f$ is coercive.
          Hence $f^\star = \ans{-2}$, reached at both minima.
  \end{itemize}
\end{frame}

\begin{frame}{\soltitle{M1}{3}{4}}
  \begin{columns}[T]
    \begin{column}{0.54\textwidth}
      \textbf{(c) One GD step from $\x_0 = (1,0)$ with $\alpha = 0.1$.}
      \[
        \g_0 = \nabla f(1,0) = (4-0,\; 0-4) = (4,-4)
      \]
      \[
        \x_1 = (1,0) - 0.1\,(4,-4) = \ans{(0.6,\;0.4)}
      \]
      \[
        f(\x_0) = 1, \qquad f(\x_1) = \ans{-0.8048}
      \]
      \footnotesize ($f(\x_1) = 0.6^4 + 0.4^4 - 4(0.24) = 0.1296 + 0.0256 - 0.96$.)
      \small We knew it would go down for small $\alpha$, because the
      directional derivative is $\nabla f\T(-\g_0) = -\norm{\g_0}^2 = -32 < 0$.
      Keep iterating with $\alpha = 0.1$ and you arrive at $(1,1)$.
    \end{column}
    \begin{column}{0.44\textwidth}
      \centering
      \includegraphics[width=0.95\linewidth]{fig_M1_step}
    \end{column}
  \end{columns}
\end{frame}

\begin{frame}{\soltitle{M1}{4}{4}}
  \textbf{(d) Local stability.} Near a minimiser, $\e_{k+1} \approx (\bI - \alpha\nabla^2 f(\x^\star))\,\e_k$,
  with eigen-factors $1 - 8\alpha$ and $1 - 16\alpha$. The stiffer one decides:
  \[
    |1-16\alpha| < 1 \;\Longleftrightarrow\; \ans{0 < \alpha < \tfrac{2}{16} = 0.125}
  \]
  The best local choice is $\alpha = \tfrac{2}{8+16} = \tfrac1{12}$, with local rate $\tfrac{16-8}{16+8} = \tfrac13$.
  \textbf{The saddle.} $\nabla f(0,0) = \mathbf 0$, so GD started there \alert{never moves at all}.
  On the line $y=-x$, $\nabla f(t,-t) = (4t^3+4t)(1,-1)$ keeps the iterates on the line, and $f(t,-t) = 2t^4+4t^2$
  is smallest at $t=0$. So GD \alert{slides into the saddle}: that line is its stable manifold.

  \begin{keybox}[Insight]
    Any nudge along $(1,1)$ grows by $1+4\alpha$ per step and escapes.
    Numerically, $(10^{-3},10^{-3})$ ends at $(1,1)$, while $(10^{-3},-10^{-3})$ ends at $(0,0)$.
    Saddles only capture a set of starting points of measure zero.
  \end{keybox}
\end{frame}
